% German -- what the reading editions' preamble leaves to the language.
% Input by lib/frank-preamble.tex as \input{\FrankLib/lang/\BookLang.tex};
% docs/languages.md section 6 lists the hooks every file here must define.
% The Latin-script model, as Italian: no font of its own, no marks to strip,
% two passes.  What is German's own is which of the two Germans babel gets.

% ---- the babel language and the switches ----------------------------------
% ngerman, not german: the two are different sets of hyphenation patterns.
% What tells them apart is ck and st.  The old spelling divided Zucker as
% Zuk-ker, a break that changes a letter and so cannot be written as a
% pattern at all, so the old set declines ck outright -- Zucker backen
% Drucker Wecker come out unbroken -- and it never divides s from t
% (Fen-ster, Mei-ster).  The reformed set, which every book printed since
% 1996 follows, gives Zu-cker ba-cken Dru-cker We-cker, Fens-ter Meis-ter.
% Schiff-fahrt is not the test it looks like: both sets break it there.
% Measured with \showhyphens under \babelprovide[import=de-1901]{german}
% against [import=de-1996]{ngerman}.
% This name is also what every \foreignlanguage and \begin{otherlanguage} in
% the core preamble is handed, so babel must know it exactly as spelt here.
\newcommand{\FrankLangName}{ngerman}
\FrankRTLfalse           % chunk left, gloss right; \pw needs no bidi isolate
\FrankHasBarefalse       % nothing is written in, so nothing can come off again
\FrankHasAltfalse        % no alternate face
\FrankHasReadingfalse
\FrankHasVertfalse

% ---- fonts ------------------------------------------------------------------
% The registry's tex.main is null for German: the text is set in the main
% roman the core preamble already loaded (TeX Gyre Pagella), the same face as
% the glosses, so a German chunk and its English gloss differ by size and
% colour only.  Pagella has the umlauts, the ß and the two letters the
% pronunciation line uses beyond them (ç and ž), so German needs no face of
% its own.  It does need the core's DejaVu Serif fallback for one letter, as
% the Italian and French lines need it for ʃ ʒ ɲ: Pagella has no capital ẞ --
% font.characters[0x1E9E] is nil for it and set for DejaVu Serif, asked of
% luaotfload on this machine.  Nothing here case-folds, so a ẞ reaches the
% page only where the source itself wrote one, a name or a title in capitals
% (STRAẞE), and there it costs the book a second embedded face.
% import=de brings babel's German hyphenation and its quotation marks.
\babelprovide[import=de]{ngerman}

% ---- the vocabulary line -----------------------------------------------------
% \vb{infinitive}{sound}{preterite}{sound}{past participle}{sound}{meaning}:
% the three principal parts, which is what a German dictionary prints and what
% every other form is built out of (geben, gab, gegeben).  The present is
% deliberately not among them: a strong verb changes its vowel there too
% (er gibt), but that form is one person of one tense, while the participle
% decides the whole perfect -- so the vowel change rides with the meaning
% instead.  docs/lang/de.md names the same three in the same order, which is
% what keeps the annotator's slots and these labels from drifting apart.
% What rides with the meaning is ONE parenthesis after it, items parted by
% "; ", each only where it is not the ordinary case:
%   er fährt            the third person present, where it is not the stem
%                       plus -t or -et (er gibt, er nimmt; ich kann for a
%                       modal), and a separable verb's whole: er fährt ab
%   aux. sein           the perfect is made with sein; aux. haben/sein where
%                       it takes both.  haben alone is never printed
%   + dat., auf + acc.  the case or the preposition the verb governs
%   \vb{fahren}{}{fuhr}{}{gefahren}{}{to drive (er fährt; aux. sein)}
% A reflexive verb carries its sich in the first slot, sich freuen, and the
% two other forms stay bare (freute, gefreut).  The chunk editor's sidebar
% drafts this same line from the dictionary (lib/verbs/de.py), in this
% order and these words, so a change to it is a change there too.
% The reader takes the same pair from the registry (lib/languages.json,
% "vb_labels": ["pret.", "p.p."], the three forms in words as "vb_forms"); the
% two must agree, or the PDF and the reader of one book label the same form
% differently.  lib/newlang.py --check compares them.
\newcommand{\FrankVbPres}{pret.}
\newcommand{\FrankVbPast}{p.p.}

% ---- the two Lua filters ---------------------------------------------------
% Both are the identity: German writes no mark the bare pass would drop (the
% registry's `strip` is null) and its labels are already in Latin digits.
% They still have to exist, because the core calls them for every language.
% Nothing here folds ß into ss, either.  That is a spelling and not a mark --
% Straße and the Swiss Strasse are both right, and a book printed before 1996
% says daß -- so a filter that tidied it would put on the page a word the
% source does not contain.  No #, % or literal backslash in a \directlua body:
% TeX reads it first (NOTES section 9 of the Persian book).
\directlua{
  function frank_strip(s) return s end
  function frank_latin(s) return s end
}

% ---- the "How to read this" page -----------------------------------------------
% Printed by lib/frank-frontmatter.tex, before the first chapter: the only
% place the reader is told what the passes are for.  German words are set in
% \textit and not in \pw, as in the Italian file: \pw shrinks a word to 8.6pt
% for a script that needs telling apart from the English, and here the two are
% the same face, where the italic is what does the telling.
\newcommand{\FrankHowTo}{%
Every passage comes twice.

\smallskip
\textbf{First} the whole sentence, and nothing else. Try to read it. Get as far
as you can and do not worry about what you cannot get — the point is the
attempt, made before any help arrives.

\smallskip
\textbf{Second} the same sentence in small chunks, German on the left, and on
the right three lines: how it is pronounced, then how to look each word up,
then what it means. Now you find out how much you had right. There is no third
pass, because nothing was written into the German for the first one and so
there is nothing to take away again: the sentence you attempted is already the
sentence as a German book prints it, capital nouns and all.

\smallskip
The pronunciation line is given only where the spelling misleads, and is empty
otherwise. It respells the chunk in small letters: a colon after a vowel means
the vowel is long (\textit{schta:t} \textit{Staat} against \textit{schtat}
\textit{Stadt}), an apostrophe stands before the stressed syllable, \textit{x}
is the \textit{ch} of \textit{Bach} and \textit{ç} the thinner one of
\textit{ich}, a final \textit{-er} is a dark vowel written \textit{a}
(\textit{'fa:ta}, \textit{Vater}). Stress alone separates
\textit{'ü:bazetsen}, to ferry someone across, from \textit{ü:ba'zetsen}, to
translate, so a verb with a prefix always gets a line.

\smallskip
Verbs are given as infinitive, preterite, past participle, meaning —
\textit{geben}, \textit{gab}, \textit{gegeben} — the three you cannot work
out from one another. What they leave out stands in brackets after the
meaning: the present where it changes its vowel too (\textit{er gibt}),
\textit{aux. sein} where the perfect takes \textit{sein} and not
\textit{haben}, and the case or the preposition the verb governs
(\textit{+ dat.}, \textit{auf + acc.}). A reflexive verb is given with its
\textit{sich}. A verb with a separable prefix is given whole,
\textit{aufstehen}, even where the sentence has pulled it in two and the
\textit{auf} waits at the end of the clause.
}
